\section*{STAR METHODS}

%%% Key Resources Table
\subsection*{Key resources table}

\begin{tabular}{|l|l|l|}
\hline
\textbf{REAGENT or RESOURCE} & \textbf{SOURCE} & \textbf{IDENTIFIER} \\
\hline
\multicolumn{3}{|l|}{\textbf{Software and algorithms}} \\
\hline
MATLAB & MathWorks & https://www.mathworks.com/products/matlab \\
\hline
Custom simulation code & This paper & https://github.com/jamesjscully/Dendronotus \\
\hline
\end{tabular}


%%% Method Details
\subsection*{Method details}

\subsubsection*{Neuron model}

The dynamics of the membrane potential, $V$, is governed by the following equation:
\begin{equation}
 	C_{m} {V}^\prime = -I_{I} - I_K - I_{T}  - I_{KCa} -  I_{h}  - I_{leak}  - I_{syn},
\end{equation}
where $C_m = 1$\,nF is the membrane capacitance.

The fast inward sodium and calcium $I_{I}$-current is given by
\begin{equation}
I_{I}=g_{I}\,h\, m^{3}_{\infty}(V)(V-E_{I}),
\end{equation}
with the reversal potential $E_{I}=30$mV and the maximal conductance value $g_{I}=4$nS, and where the infinite functions are given by
\begin{align}
 m_{\infty}(V) &= \frac{\alpha_{m}(V)}{\alpha_{m}(V) +\beta_{m}(V)},   \\
 \alpha_{m}(V) &= 0.1 \frac{50-V_s}{-1+e^{(50-V_s)/10}}, \\
 \beta_{m}(V) &= 4 e^{(25-V_s)/18}.
\end{align}

The dynamics of its inactivation gating variable $h$ is given by
\begin{align}
 {h}^\prime &=  \frac{h_{\infty}(V)-h}{ {\tau_{h}(V)} }, \\
h_{\infty}(V) &= \frac{\alpha_{h}(V)}{\alpha_{h}(V) +\beta_{h}(V)}, \\
 \tau_{h}(V) &= \frac{12.5}{\alpha_{h}(V) +\beta_{h}(V)},
\end{align}
where
\begin{align}
 \alpha_{h}(V) &= 0.07 e^{(25-V_s)/20}, \\
 \beta_{h}(V) &= \frac{1}{1+e^{(55-V_s)/10}}, \\
 V_s &= \frac{127V+8265}{105}{\rm mV}.
\end{align}

The fast potassium $I_K$-current is given by the equation
\begin{equation}
I_K=g_{K} n^{4}(V-E_{K}),
\end{equation}
with the reversal potential $E_{K}=-75$mV and the maximal conductance set as $g_K=0.3$nS.
The dynamics of inactivation gating variable are described by
\begin{align}
{n}^\prime &= \frac{n_{\infty}(V)-n}{\tau_{n}(V)}, \\
n_{\infty}(V) &= \frac{\alpha_{n}(V)}{\alpha_{n}(V) +\beta_{n}(V)}, \\
\tau_{n}(V) &= \frac{12.5}{\alpha_{n}(V) +\beta_{n}(V)}, \\
\alpha_{n}(V) &= 0.01 \frac{55-V_s}{e^{(55-V_s)/10}-1}, \\
\beta_{n}(V) &= 0.125 e^{(45-V_s)/80}.
\end{align}

The rapidly depolarizing h-current is given by
\begin{equation}
I_{h}=g_{h} \frac{y (V-E_{h})}{(1+e^{-(V-63)/7.8})^3},
\end{equation}
with $E_{h}=70$mV, and $g_h=0.0006$nS; the dynamics of its $y$-activation probability is described by
\begin{equation}
{y}^\prime = \frac{1}{2} \left [ \frac{1}{1+e^{10(V-50)}}-y \right ]/ \left [ 7.1+\frac{10.4}{1+e^{(V+68)/2.2}} \right ].
\end{equation}

The leak current is given by
\begin{equation}
I_{leak} =g_{L} (V-E_{L}),
\end{equation}
with $E_{L} =-40{\rm mV}$, and $g_{L}=0.003{\rm nS}$.

The sub-group of slow currents in the model includes the TTX-resistant sodium and calcium $I_{T}$-current given by
\begin{equation}
 I_{T} = g_{T} x (V-E_{I}),
\end{equation}
with $E_{I} =30{\rm mV}$ and $g_{T}=0.01{\rm nS}$, where the dynamics of its slow activation variable are described by
\begin{align}
{x}^\prime &= \frac{x_{\infty}(V)-x}{\tau_{x}}, \\
x_{\infty} &= \frac{1}{1+e^{-0.15(V+50-\Delta_x)}},
\end{align}
with the time constant $\tau_{x} = 100$\,ms.

The slowest outward ${\rm Ca}^{2+}$ activated $K^+$ current given by
\begin{equation}
I_{KCa} =g_{KCa}\frac{[\rm Ca]_i}{0.5+[\rm Ca]_i}(V-E_{K}),
\end{equation}
with $E_{K}=-75\mbox{mV}$, and $g_{KCa}=0.03$, where the dynamics of intracellular calcium concentration is governed by
\begin{equation}
{[\rm Ca]}^\prime  =  \rho \left ( K_{c}\, x\, (E_{\rm Ca}-V + \Delta_{\rm Ca})-[\rm Ca] \right  )
\end{equation}
with the Nernst reversal potential $E_{Ca}=140$mV, and small constants $\rho= 0.0003{\rm ms}^{-1}$ and $K_c=0.0085{\rm mV}^{-1}$.


\subsubsection*{Synapse models}

\paragraph{Fast $\alpha$-synapses.}
Fast excitatory and inhibitory synapses are modeled using the $\alpha$-synapse formalism \cite{destexhe1994synthesis, destexhe1998kinetic, Rall67}.
The synaptic current is given by
\begin{equation}
I_{syn} = g_{syn} \cdot s \cdot (V_{post} - E_{syn}),
\end{equation}
where $g_{syn}$ is the maximal synaptic conductance, $s$ is the synaptic gating variable, and $E_{syn}$ is the reversal potential ($E_{syn} = +40$mV for excitation, $E_{syn} = -80$mV for inhibition).

The synaptic gating variable $s$ follows
\begin{align}
{s}^\prime &= \alpha \cdot f_{\infty}(V_{pre}) \cdot (1 - s) - \beta \cdot s,
\end{align}
where $f_{\infty}(V_{pre})$ is a sigmoid function of presynaptic voltage:
\begin{equation}
f_{\infty}(V_{pre}) = \frac{1}{1 + e^{-k(V_{pre} - V_{th})}},
\end{equation}
with threshold $V_{th} = -20$\,mV and slope $k = 20$\,mV$^{-1}$.

\paragraph{Logistic synapses.}
The contralateral E/I module synapses (both Si3$\to$Si2 excitation and Si2$\to$Si3 inhibition) use a novel logistic synapse model that captures synaptic accumulation dynamics:
\begin{equation}
{s}^\prime = \alpha \cdot s \cdot (1-s) \cdot f_{\infty}(V_{pre}) - \beta \cdot (s - s_0),
\end{equation}
where $f_{\infty}(V_{pre})$ is the same sigmoid function of presynaptic voltage, and $s_0 \ll 1$ is a baseline release rate.
The key difference from the standard $\alpha$-synapse is the $s(1-s)$ term, which creates logistic growth dynamics: the synapse accumulates slowly when $s$ is small, reaches maximal accumulation rate at $s = 0.5$, and saturates as $s$ approaches 1.
The Si2$\to$Si3 inhibitory synapses use slower kinetics ($\alpha = 0.012$\,ms$^{-1}$) that produce summation on timescales of hundreds of milliseconds, setting the burst period.
The Si3$\to$Si2 excitatory synapses use faster kinetics ($\alpha = 0.05$\,ms$^{-1}$) but retain the logistic accumulation dynamics for consistent E/I module behavior.

\paragraph{Synaptic normalization.}
Without normalization, the synaptic gating variable $s$ is bounded between 0 and a maximum value that depends on the kinetic parameters $\alpha$ and $\beta$.
Since the synaptic conductance is proportional to $g_{syn} \cdot s$, changing $\alpha$ and $\beta$ would inadvertently affect the effective synaptic strength, creating unwanted codependencies between kinetic timescales and synaptic amplitude.
To decouple synaptic strength from kinetic parameters, we normalize the synaptic current:
\begin{equation}
I_{syn} = g_{syn} \cdot \frac{s}{\text{scale}_s} \cdot (V_{post} - E_{syn}),
\end{equation}
where the scaling factors depend on synapse type:
\begin{align}
\text{scale}_{\text{fast}} &= \frac{\alpha}{\alpha + \beta} \quad \text{(for $\alpha$-synapses)}, \\
\text{scale}_{\text{slow}} &= \frac{\alpha - \beta}{\alpha} \quad \text{(for logistic synapses)}.
\end{align}
These factors normalize the theoretical asymptotic limit to unity for all parameter combinations, ensuring that $g_{syn}$ directly represents the maximum synaptic conductance independent of the kinetic parameters.


\subsubsection*{Network connectivity}

The \textit{Dendronotus iris} swim CPG model consists of four core interneurons (Si2L, Si2R, Si3L, Si3R) with the following synaptic connections:
\begin{itemize}
    \item \textbf{Contralateral excitation}: Si3R$\to$Si2L and Si3L$\to$Si2R (logistic synapses)
    \item \textbf{Contralateral inhibition}: Si2L$\to$Si3R and Si2R$\to$Si3L (logistic synapses)
    \item \textbf{Reciprocal Si3 inhibition}: Si3L$\longleftrightarrow$Si3R ($\alpha$-synapses)
    \item \textbf{Reciprocal Si2 inhibition}: Si2L$\longleftrightarrow$Si2R ($\alpha$-synapses)
    \item \textbf{Electrical coupling}: Si3R$\longleftrightarrow$Si2L and Si3L$\longleftrightarrow$Si2R (gap junctions)
\end{itemize}

The command neuron Si1 provides neuromodulatory input that scales the effective strength of Si3$\to$Si2 excitation, enabling frequency control.


\subsubsection*{Network assembly and calibration}

We assembled the four-cell swim CPG network by integrating validated two-cell subnetworks---specifically half-center oscillators (HCOs) and excitatory-inhibitory (E/I) modules---following the approach in \cite{scully2024pairing}.
Synaptic kinetics for mutual inhibition and E/I interactions were matched to previously characterized modules.

Cellular parameters were preserved from prior models, with the exception of Si3 neurons, whose endogenous excitability was increased by shifting the voltage threshold $\Delta[\text{Ca}]$ by 5\,mV to promote spiking.

\paragraph{Almost-oscillatory sub-circuits.}
By ``almost-oscillatory'' sub-circuits, we mean configurations that exhibit damped oscillatory behavior or subthreshold rhythmic fluctuations when isolated, but do not sustain stable limit-cycle oscillations without network interactions.
These sub-circuits lie close to oscillatory regimes in parameter space, near bifurcations that separate quiescent from rhythmic states, such that modest network coupling can recruit them into coordinated bursting.
This design principle contrasts with CPG architectures built from intrinsically bursting pacemaker neurons that generate robust oscillations independently of network context.

\paragraph{Calibration protocol.}
Network calibration was performed by symmetrically tuning synaptic conductances using the following protocol:
\begin{enumerate}
    \item \textbf{Configuration:} Pre- and postsynaptic voltage variables were assigned according to the known swim circuits of \textit{Dendronotus iris} \cite{SakuraiKatz2016}.
    \item \textbf{Initialization:} All cells were set to a quiescent state with synaptic conductances at zero.
    \item \textbf{Drive:} Synaptic drive was applied to Si3 until the excitatory postsynaptic gating variable $s$ showed sustained accumulation, measured by local minima rising above 0.1 and saturating over 0.5--3\,s.
    \item \textbf{Excitation:} Excitatory conductance from Si3 to contralateral Si2 was increased until inhibitory $s$ (Si2-to-Si2) reached minima above 0.1 over 0.5--3\,s.
    \item \textbf{Slow inhibition:} Inhibitory conductance in the E/I pair was raised until Si3 frequency decreased by 50\% over 1\,s.
    \item \textbf{Fast inhibition:} Inhibitory conductances for both HCOs (Si2L--Si2R and Si3L--Si3R) were elevated until a stable burst rhythm emerged, visually matching empirical voltage traces \cite{SakuraiKatz2016}.
    \item \textbf{Fine-tuning:} Excitability was adjusted slightly (e.g., $\Delta[\text{Ca}] < 2$\,mV) to achieve a duty cycle matching experimental observations.
\end{enumerate}


\subsubsection*{Model validation}

To validate the CPG model, we performed direct current injection into the voltage equation of selected interneurons, adding a current term to the right-hand side of the membrane potential equation and varying its amplitude to mimic experimental perturbations.
No additional parameter tuning, optimization, or fitting procedures were employed during these tests.
This approach ensures that the model's responses reflect its intrinsic dynamics and structure, maintaining the integrity of the validation process.

All validation responses shown are typical: simulations use random initial conditions for the four interneurons and synapses, plus calibrated colored noise (up to 20 Hz), to probe robustness.
This explains trial-to-trial variability in initial phase and demonstrates that the model's behavior is not dependent on specific initial conditions.

Validation focuses on qualitative reproduction of experimental phenomena through systematic comparison of model output with published voltage traces and perturbation responses \cite{SakuraiKatz2016, sakurai2019command, sakurai2022bursting}.
Quantitative parameter fitting was not attempted due to limited availability of quantitative measurements in the experimental literature.
Our approach prioritizes mechanistic insight---demonstrating that the proposed circuit architecture and synaptic dynamics can reproduce the diversity of observed behaviors---over numerical precision in matching specific parameter values.


\subsubsection*{Neuromodulation}

Si1 neurons modulate the Si3$\to$Si2 excitatory synapses through a dynamic modulatory variable $s_m \in [0, 1]$ that evolves according to
\begin{equation}
{s_m}^\prime = \alpha_m \cdot s_m \cdot (1 - s_m) \cdot f_{\infty}(V_{Si1}) - \beta_m \cdot (s_m - s_{m}^0),
\end{equation}
where $f_{\infty}(V_{Si1})$ is the sigmoid function of the Si1 presynaptic voltage, and $0 < s_m^0 \ll 1$ is a baseline value.

In the MATLAB implementation, modulation is normalized by
\begin{equation}
\text{scale}_m = \frac{\alpha_m - \beta_m}{\alpha_m},
\end{equation}
and the gain parameter $g_m$ enters through the factor $(1 + g_m s_m / \text{scale}_m)$.
When $s_m = 0$, the synapse operates at baseline strength.
We compare two modulatory mechanisms that target distinct biophysical processes:

\paragraph{Postsynaptic modulation.}
The modulatory variable scales the maximal conductance of the logistic Si3$\to$Si2 synapse directly:
\begin{equation}
I_{syn} = g_{syn} \cdot \left(1 + g_m \frac{s_m}{\text{scale}_m}\right) \cdot \frac{s}{\text{scale}_s} \cdot (V_{post} - E_{syn}),
\end{equation}
We term this ``postsynaptic'' because the synaptic conductance $g_{syn}$ characterizes the electrical properties of the postsynaptic membrane---specifically, the density and conductance of postsynaptic receptors and ion channels that determine the amplitude of the postsynaptic response.

\paragraph{Presynaptic modulation.}
Alternatively, the modulatory variable scales the accumulation rate $\alpha$ of the same logistic synapse:
\begin{equation}
{s}^\prime = \alpha \cdot \left(1 + g_m \frac{s_m}{\text{scale}_m}\right) \cdot s \cdot (1-s) \cdot f_{\infty}(V_{pre}) - \beta \cdot (s - s_0).
\end{equation}
We term this ``presynaptic'' because the gating variable $s(t)$ models the dynamics of neurotransmitter release, a process that occurs at the presynaptic terminal.
Modulating the $\alpha$ parameter affects the rate at which neurotransmitter accumulates, mimicking biological mechanisms such as changes in vesicle priming, calcium sensitivity, or release probability.

Both mechanisms increase effective synaptic strength, but presynaptic modulation additionally accelerates synaptic accumulation, producing qualitatively different effects on network dynamics as explored in the Results.

\subsubsection*{Neuromodulation scan and rate-matching procedure}

For the neuromodulation comparison in Fig.~\ref{fig:si1_neuromod}, we analyzed the circuit under sustained Si1 drive and compared two modulation targets within the Si3$\to$Si2 excitatory pathway.
In the presynaptic condition, the gain parameter multiplied the logistic accumulation-rate parameter $\alpha$.
In the postsynaptic condition, the same gain parameter multiplied the effective synaptic conductance.
For each condition, we scanned gain from 0 to 8 and measured burst frequency from successive Si2 burst intervals after excluding early transients.

To compare model output with experiment, we used the Si1 and Si2 traces digitized from preparation 130618-03 together with the published Si1-versus-burst summary data from Sakurai \& Katz (2019).
Burst onsets were detected from Si2 threshold crossings, and Si1 firing rate was computed within each inter-burst interval.
This produced paired measurements of Si1 spike frequency and burst frequency for both the biological traces and the representative model traces shown in Fig.~\ref{fig:si1_neuromod}E.
The same detection procedure was applied to both model variants.
For display and interval matching, the biological trace time base was stretched by a factor of 2 relative to the raw digitized recording; accordingly, frequencies extracted from the published summary points were divided by 2 so that the trace-based and summary-based biological comparisons were expressed on the same timescale.

\paragraph{Representative traces.}
The representative presynaptic and postsynaptic traces in Fig.~\ref{fig:si1_neuromod}E were generated by replaying the same digitized Si1 voltage trace into both model variants.
The representative presynaptic trace used control gain $g_m = 3.545$, chosen so that the modulated excitatory accumulation rate reached the cited target value under sustained drive.
The representative postsynaptic trace used control gain $g_m = 3.0$, together with conductance compensation in the unmodulated baseline, so that the effective excitatory conductance matched the cited target under the same saturated-drive assumption.

\paragraph{Figure-specific parameters.}
The clean neuromodulation dashboard figure was generated from the calibrated Julia analysis pipeline using the saved parameter set listed below.
These values define the exact scans and representative traces used for Fig.~\ref{fig:si1_neuromod} and should be taken as the figure-specific parameter source of truth for this analysis.

\begin{table}[htbp]
\centering
\caption{Figure 5 neuromodulation scan parameters}
\label{tab:neuromod_dashboard_params}
\begin{tabular}{|l|l|l|}
\hline
\textbf{Parameter} & \textbf{Value} & \textbf{Role in analysis} \\
\hline
$\Delta x_{\mathrm{Si2}}$ & $-4.0$ mV & Si2 slow-current voltage shift \\
$\Delta x_{\mathrm{Si3}}$ & $-4.0$ mV & Si3 slow-current voltage shift \\
$g_{0,\mathrm{pre}}$ & 0.0008 & Si1$\to$Si3 drive, presynaptic condition \\
$g_{0,\mathrm{post}}$ & 0.0016 & Si1$\to$Si3 drive, postsynaptic condition \\
$g_{m,\mathrm{pre}}^{\ast}$ & 3.545 & representative presynaptic trace gain \\
$g_{m,\mathrm{post}}^{\ast}$ & 3.0 & representative postsynaptic trace gain \\
$\alpha_1$ & 0.008 ms$^{-1}$ & Si1 synaptic activation rate \\
$\beta_1$ & 0.0016 ms$^{-1}$ & Si1 synaptic decay rate \\
$s_0^{\mathrm{floor}}$ & 0.013 & baseline Si1 synaptic floor \\
$\alpha_m$ & 0.006 ms$^{-1}$ & modulatory buildup rate \\
$\beta_m$ & $5.0\times10^{-5}$ ms$^{-1}$ & modulatory decay rate \\
$s_m^{\mathrm{floor}}$ & 0.0213 & baseline modulatory floor \\
$\Delta\mathrm{Ca}_{\mathrm{Si2}}$ & $-25.0$ mV & calcium shift, Si2 pair \\
$\Delta\mathrm{Ca}_{\mathrm{Si3}}$ & $-25.0$ mV & calcium shift, Si3 pair \\
$g_{\mathrm{base}}^{\mathrm{pre}}$ & 0.001 & baseline Si3$\to$Si2 conductance, presynaptic condition \\
$g_{\mathrm{base}}^{\mathrm{post}}$ & 0.04 & baseline Si3$\to$Si2 conductance, postsynaptic condition \\
$\alpha_{x,\mathrm{pre}}$ & 0.011 ms$^{-1}$ & baseline logistic accumulation, presynaptic condition \\
$\alpha_{x,\mathrm{post}}$ & 0.011 ms$^{-1}$ & baseline logistic accumulation, postsynaptic condition \\
$\beta_{x,\mathrm{pre}}$ & 0.001 ms$^{-1}$ & logistic decay, presynaptic condition \\
$\beta_{x,\mathrm{post}}$ & 0.001 ms$^{-1}$ & logistic decay, postsynaptic condition \\
$g_{\mathrm{inh,slow}}$ & 0.005 & Si2$\to$Si3 slow inhibition \\
$\alpha_{\mathrm{inh,slow}}$ & 0.015 ms$^{-1}$ & Si2$\to$Si3 inhibitory accumulation \\
$\beta_{\mathrm{inh,slow}}$ & 0.0005 ms$^{-1}$ & Si2$\to$Si3 inhibitory decay \\
$g_{\mathrm{Si2\leftrightarrow Si2}}$ & 0.0115 & reciprocal Si2 inhibition \\
$\alpha_{\mathrm{Si2\leftrightarrow Si2}}$ & 0.01 ms$^{-1}$ & reciprocal Si2 inhibitory rise \\
$\beta_{\mathrm{Si2\leftrightarrow Si2}}$ & 0.01 ms$^{-1}$ & reciprocal Si2 inhibitory decay \\
$g_{\mathrm{Si3\leftrightarrow Si3}}$ & 0.005 & reciprocal Si3 inhibition \\
$\alpha_{\mathrm{Si3\leftrightarrow Si3}}$ & 0.01 ms$^{-1}$ & reciprocal Si3 inhibitory rise \\
$\beta_{\mathrm{Si3\leftrightarrow Si3}}$ & 0.002 ms$^{-1}$ & reciprocal Si3 inhibitory decay \\
$t_1$ & 1000 ms & synaptic timescale parameter in the scan model \\
\hline
\end{tabular}
\end{table}


\subsubsection*{Simulation and analysis tools}

The core circuit-development and validation simulations were performed using custom MATLAB code.
For the neuromodulation dashboard analysis in Fig.~\ref{fig:si1_neuromod}, we used a Julia implementation of the same model equations to run parameter scans, generate representative traces, and extract matched Si1-rate versus burst-rate measurements.
Differential equations were integrated numerically with fixed trace sampling for visualization and burst detection.
Burst detection and period measurement were performed using threshold-crossing algorithms on the membrane voltage traces.


%%% Model Parameters
\subsection*{Model parameters}

Table~\ref{tab:intrinsic_params} lists intrinsic neuronal parameters, which are identical across all four interneurons except where noted.
Table~\ref{tab:synaptic_params} lists synaptic parameters for network connectivity.

\begin{table}[htbp]
\centering
\caption{Intrinsic neuronal parameters}
\label{tab:intrinsic_params}
\begin{tabular}{|l|l|l|l|}
\hline
\textbf{Parameter} & \textbf{Description} & \textbf{Value} & \textbf{Units} \\
\hline
\multicolumn{4}{|l|}{\textit{Fast inward current ($I_I$)}} \\
\hline
$g_I$ & Maximal conductance & 4 & nS \\
$E_I$ & Reversal potential & 30 & mV \\
\hline
\multicolumn{4}{|l|}{\textit{Fast potassium current ($I_K$)}} \\
\hline
$g_K$ & Maximal conductance & 0.3 & nS \\
$E_K$ & Reversal potential & $-75$ & mV \\
\hline
\multicolumn{4}{|l|}{\textit{Leak current ($I_{leak}$)}} \\
\hline
$g_L$ & Maximal conductance & 0.003 & nS \\
$E_L$ & Reversal potential & $-40$ & mV \\
\hline
\multicolumn{4}{|l|}{\textit{TTX-resistant current ($I_T$)}} \\
\hline
$g_T$ & Maximal conductance & 0.01 & nS \\
$\tau_x$ & Activation time constant & 100 & ms \\
$\Delta_x$ & Voltage shift & $-4$ & mV \\
\hline
\multicolumn{4}{|l|}{\textit{Calcium-activated potassium current ($I_{KCa}$)}} \\
\hline
$g_{KCa}$ & Maximal conductance & 0.03 & nS \\
$\rho$ & Calcium dynamics rate & 0.0003 & ms$^{-1}$ \\
$K_c$ & Calcium scaling & 0.0085 & mV$^{-1}$ \\
$E_{Ca}$ & Calcium reversal potential & 140 & mV \\
\hline
\multicolumn{4}{|l|}{\textit{h-current ($I_h$)}} \\
\hline
$g_h$ & Maximal conductance & 0.0006 & nS \\
$E_h$ & Reversal potential & 70 & mV \\
\hline
\multicolumn{4}{|l|}{\textit{Calcium shift ($\Delta_{Ca}$) --- sets intrinsic excitability}} \\
\hline
 & Si3L, Si3R & $-55$, $-60$ & mV \\
 & Si2L, Si2R & $-30$, $-22$ & mV \\
\hline
\end{tabular}
\end{table}

\begin{table}[htbp]
\centering
\caption{Synaptic and network parameters}
\label{tab:synaptic_params}
\begin{tabular}{|l|l|l|l|l|}
\hline
\textbf{Synapse} & \textbf{Type} & $\boldsymbol{\alpha}$ \textbf{(ms$^{-1}$)} & $\boldsymbol{\beta}$ \textbf{(ms$^{-1}$)} & $\boldsymbol{g_{syn}}$ \textbf{(nS)} \\
\hline
\multicolumn{5}{|l|}{\textit{Contralateral E/I modules}} \\
\hline
Si3R$\to$Si2L & Logistic (exc) & 0.05 & 0.001 & 0.005 \\
Si3L$\to$Si2R & Logistic (exc) & 0.05 & 0.001 & 0.005 \\
Si2L$\to$Si3R & Logistic (inh) & 0.012 & 0.001 & 0.01 \\
Si2R$\to$Si3L & Logistic (inh) & 0.012 & 0.001 & 0.01 \\
\hline
\multicolumn{5}{|l|}{\textit{Reciprocal inhibition (half-center oscillators)}} \\
\hline
Si3L$\leftrightarrow$Si3R & $\alpha$-synapse (inh) & 0.01 & 0.004 & 0.02 \\
Si2L$\leftrightarrow$Si2R & $\alpha$-synapse (inh) & 0.01 & 0.005 & 0.02 \\
\hline
\multicolumn{5}{|l|}{\textit{Electrical coupling (gap junctions)}} \\
\hline
Si3R$\leftrightarrow$Si2L & Gap junction & \multicolumn{2}{l|}{---} & 0.002 \\
Si3L$\leftrightarrow$Si2R & Gap junction & \multicolumn{2}{l|}{---} & 0.002 \\
\hline
\end{tabular}
\vspace{0.5em}

\small
Reversal potentials: $E_{exc} = +40$\,mV (excitatory), $E_{inh} = -80$\,mV (inhibitory).
Logistic synapses use $s_0 = 0.0001$ as baseline.
Sigmoid function: $f_{\infty}(V) = 1/(1 + e^{-k(V - V_{th})})$ with $V_{th} = -20$\,mV and slope $k = 20$\,mV$^{-1}$.
\end{table}


%%% Quantification and Statistical Analysis
\subsection*{Quantification and statistical analysis}

Burst periods were computed as the time interval between successive burst onsets in the Si3 neurons.
Mean periods and variability were calculated over multiple burst cycles after discarding initial transients.
